% =============================================================================
% Wstęp — sekcja do wklejenia w istniejący dokument LaTeX
%
% W dokumencie nadrzędnym użyj jednego z poniższych wariantów:
%   \chapter{Wstęp}   % klasa report / book
%   \section{Wstęp}   % klasa article
%
% Wymagane w preambule dokumentu nadrzędnego:
%   \usepackage[polish]{babel}
%   \usepackage{amsmath, amssymb}
%
% Dołączenie: % =============================================================================
% Wstęp — sekcja do wklejenia w istniejący dokument LaTeX
%
% W dokumencie nadrzędnym użyj jednego z poniższych wariantów:
%   \chapter{Wstęp}   % klasa report / book
%   \section{Wstęp}   % klasa article
%
% Wymagane w preambule dokumentu nadrzędnego:
%   \usepackage[polish]{babel}
%   \usepackage{amsmath, amssymb}
%
% Dołączenie: % =============================================================================
% Wstęp — sekcja do wklejenia w istniejący dokument LaTeX
%
% W dokumencie nadrzędnym użyj jednego z poniższych wariantów:
%   \chapter{Wstęp}   % klasa report / book
%   \section{Wstęp}   % klasa article
%
% Wymagane w preambule dokumentu nadrzędnego:
%   \usepackage[polish]{babel}
%   \usepackage{amsmath, amssymb}
%
% Dołączenie: % =============================================================================
% Wstęp — sekcja do wklejenia w istniejący dokument LaTeX
%
% W dokumencie nadrzędnym użyj jednego z poniższych wariantów:
%   \chapter{Wstęp}   % klasa report / book
%   \section{Wstęp}   % klasa article
%
% Wymagane w preambule dokumentu nadrzędnego:
%   \usepackage[polish]{babel}
%   \usepackage{amsmath, amssymb}
%
% Dołączenie: \input{wstep}
% =============================================================================

% --- Odkomentuj odpowiednią linię lub umieść treść poniżej pod własnym nagłówkiem ---

% \chapter{Wstęp}
% \section{Wstęp}

Rynek walutowy stanowi jeden z najbardziej płynnych i globalnie zintegrowanych segmentów rynków finansowych.
Kurs wymiany dolara amerykańskiego względem złotego polskiego (USD/PLN) ma bezpośrednie znaczenie dla polskiej gospodarki:
wpływa na koszty importu, wartość zobowiązań zadłużonych w walucie obcej oraz poziom inflacji importowanej.
Instytucje finansowe, przedsiębiorstwa i inwestorzy indywidualni poszukują narzędzi wspierających podejmowanie decyzji na tym rynku,
jednak prognozowanie kursów walutowych pozostaje zadaniem wyjątkowo trudnym.

Szeregi czasowe notowań finansowych charakteryzują się wysokim poziomem szumu, nieliniowością zależności
oraz okresową zmianą reżimów rynkowych. Klasyczne modele ekonometryczne, oparte na założeniach stacjonarności
i liniowości, często nie wystarczają do uchwycenia złożonych wzorców zachowania cen.
W ostatnich latach coraz większą popularnością cieszą się metody uczenia głębokiego, w szczególności sieci rekurencyjne
typu LSTM (Long Short-Term Memory) oraz architektury oparte na mechanizmie uwagi (attention), takie jak Transformer.
Pozwalają one na modelowanie wielowymiarowych sekwencji czasowych i łączenie informacji z różnych horyzontów temporalnych.

\section*{Problem badawczy}

Niniejsza praca podejmuje problem prognozowania kierunku zmiany kursu USD/PLN w horyzoncie 14~dni roboczych.
Główne pytanie badawcze brzmi: \emph{czy na podstawie 30-dniowej historii cech technicznych
i makroekonomicznych można skutecznie przewidzieć, czy kurs wzrośnie, czy spadnie w ciągu kolejnych 14~dni?}

Punktem wyjścia jest skumulowany zwrot procentowy kursu zamknięcia w horyzoncie prognozy.
Ze względu na praktyczną interpretację wyników (sygnały kupna i sprzedaży) zadanie sprowadzono
do klasyfikacji binarnej: klasa~1 oznacza wzrost kursu, klasa~0 oznacza spadek.
Etykieta dla dnia $t$ definiowana jest jako:
\begin{equation}
y_t =
\begin{cases}
1 & \text{gdy } \dfrac{P_{t+14}}{P_t} - 1 > 0, \\[6pt]
0 & \text{gdy } \dfrac{P_{t+14}}{P_t} - 1 \le 0.
\end{cases}
\label{eq:target}
\end{equation}
gdzie $P_t$ oznacza kurs zamknięcia pary USD/PLN w dniu~$t$.

Skumulowany zwrot w procentach wyraża się natomiast wzorem:
\begin{equation}
    r_t = \left(\frac{P_{t+14}}{P_t} - 1\right) \cdot 100\,\%.
    \label{eq:return}
\end{equation}

\section*{Cel pracy}

Celem pracy jest zaprojektowanie, implementacja i ocena hybrydowego modelu głębokiego łączącego warstwę LSTM
z blokiem Transformer w zadaniu klasyfikacji kierunku kursu USD/PLN.
W ramach realizacji celu:
\begin{itemize}
    \item przygotowano wielowymiarowy zestaw cech obejmujący wskaźniki analizy technicznej
          oraz stopy procentowe banków centralnych (Fed, EBC, NBP),
    \item zbudowano i wytrenowano model na historycznych danych dziennych,
    \item oceniono skuteczność predykcji za pomocą metryk klasyfikacji
          (accuracy, precision, recall, F1) oraz porównano wyniki z prostymi benchmarkami
          (strategia ,,zawsze wzrost'' i ,,zawsze spadek'').
\end{itemize}

\section*{Opis metody}

\subsection*{Dane wejściowe}

Analizie poddano dzienne notowania z pliku \texttt{dane\_ekonomiczne.csv}, którego źródłem jest Yahoo Finance oraz World Bank Data.
Dane obejmują m.in.\ ceny OHLC pary USD/PLN, skorelowane pary walutowe (EUR/PLN, EUR/USD)
oraz stopy procentowe banków centralnych Stanów Zjednoczonych, strefy euro i Polski.
Z zestawu usunięto obserwacje przypadające na weekendy, aby zachować ciągłość danych sesyjnych.
Przed przekazaniem do modelu wszystkie cechy numeryczne standaryzowano za pomocą \texttt{StandardScaler},
dopasowanego wyłącznie na zbiorze treningowym.

\subsection*{Inżynieria cech}

Dla każdego dnia obliczono łącznie 24~zmienne wejściowe, w tym:
\begin{itemize}
    \item dzienne stopy zwrotu kursów zamknięcia (USD/PLN, EUR/PLN, EUR/USD),
    \item wskaźniki momentum za okresy 3, 5 i 10~dni,
    \item RSI (14~okresów), ATR (14~okresów), pozycję i szerokość pasm Bollingera (okno 20~dni),
    \item wskaźniki MACD wraz z linią sygnałową i histogramem,
    \item zmienność kroczącą, $z$-score (okno 20~dni), dzienny zakres High/Low oraz stosunek High/Low,
    \item stopy procentowe: US\_FED, EA\_ECB, PL\_NBP,
    \item opóźnione kierunki zmian ceny (lagi 1, 2, 3 i 5~dni).
\end{itemize}

\subsection*{Architektura modelu}

Model przyjmuje na wejściu sekwencję $X_t \in \mathbb{R}^{30 \times 24}$,
tj.\ 30~kolejnych dni obserwacji po 24~cechy (okno przesuwne).
Architektura hybrydowa składa się z następujących etapów:
\begin{enumerate}
    \item warstwa LSTM z 128~jednostkami ukrytymi (\texttt{return\_sequences=True}),
          wychwytująca lokalne zależności czasowe w sekwencji,
    \item warstwa \texttt{Dropout} ($p = 0{,}3$) redukująca przeuczenie,
    \item kodowanie pozycyjne (\emph{Positional Encoding}) nadające informację o położeniu obserwacji w oknie,
    \item blok Transformer z 8~głowicami uwagi (\emph{Multi-Head Attention}),
          $d_{\text{model}} = 128$ i warstwą sprzężeń zwrotnych o wymiarze 256,
          modelujący długoterminowe zależności w obrębie 30-dniowego okna,
    \item agregacja \texttt{GlobalAveragePooling1D}, po której następują warstwy gęste
          (64 i 16~neuronów, aktywacja ReLU, regularyzacja L2) oraz warstwa wyjściowa
          z aktywacją sigmoid (jedna neurona, klasyfikacja binarna).
\end{enumerate}

Połączenie LSTM i Transformera uzasadnia fakt, że sieć rekurencyjna efektywnie koduje krótkoterminowe wzorce cenowe,
natomiast mechanizm self-attention pozwala ważyć znaczenie poszczególnych dni w historycznym oknie,
co jest istotne przy analizie wielotygodniowych trendów.

\subsection*{Trenowanie i ewaluacja}

Dane podzielono chronologicznie: 80\,\% stanowi zbiór treningowy, 20\,\%~--- testowy.
Z części treningowej wydzielono 10\,\% obserwacji na walidację.
Ze względu na naturę szeregu czasowego podczas uczenia nie stosowano losowego mieszania próbek (\texttt{shuffle=False}).

Model trenowano algorytmem Adam z początkową szybkością uczenia $5 \times 10^{-5}$,
funkcją straty \texttt{binary\_crossentropy} oraz wagami klas (\texttt{compute\_class\_weight})
kompensującymi niezbalansowanie etykiet.
Zastosowano wczesne zatrzymanie (\emph{early stopping}) monitorujące dokładność walidacyjną
oraz redukcję szybkości uczenia przy stagnacji straty.
Predykcje ciągłe z warstwy sigmoid konwertowano na klasy za pomocą progu $0{,}5$.

Jakość modelu oceniano metrykami: accuracy, precision, recall, F1 oraz macierzą pomyłek.
Wyniki porównywano z naiwnymi strategiami bazowymi przewidującymi zawsze wzrost lub zawsze spadek kursu.

\section*{Zakres i ograniczenia}

Praca koncentruje się wyłącznie na prognozie \emph{kierunku} zmiany ceny, a nie na estymacji jej poziomu
ani wielkości zwrotu. Nie uwzględniono kosztów transakcyjnych, spreadu bid-ask ani dźwigni finansowej,
co ogranicza bezpośrednią aplikowalność wyników w realnym handlu algorytmicznym.

Model wytrenowano na danych historycznych z okresu od 2020~roku; jego skuteczność w przyszłości
nie jest gwarantowana ze względu na ryzyko przeuczenia oraz możliwość zmiany reżimu rynkowego
(np.\ interwencje banków centralnych, kryzysy geopolityczne). Wyniki należy traktować jako eksperyment badawczy,
a nie rekomendację inwestycyjną.

\section*{Struktura pracy}

% --- Dostosuj nazwy rozdziałów do własnego spisu treści ---

Praca składa się z następujących rozdziałów:
\begin{enumerate}
    \item Wstęp --- motywacja, problem badawczy, cel i zakres pracy.
    \item Przegląd literatury i podstawy teoretyczne --- szeregi czasowe finansowe,
          sieci LSTM, architektura Transformer oraz wybrane wskaźniki analizy technicznej.
    \item Przygotowanie danych i inżynieria cech --- opis źródeł danych, preprocessingu
          i konstrukcji 24~zmiennych wejściowych.
    \item Architektura i trenowanie modelu --- szczegółowy opis hybrydy LSTM\,+\,Transformer,
          hiperparametrów i procedury walidacji.
    \item Wyniki i analiza --- metryki klasyfikacji, macierz pomyłek, porównanie z benchmarkami
          oraz wizualizacja sygnałów kupna i sprzedaży.
    \item Podsumowanie i wnioski --- ocena osiągniętych rezultatów i kierunki dalszych badań.
\end{enumerate}

% =============================================================================

% --- Odkomentuj odpowiednią linię lub umieść treść poniżej pod własnym nagłówkiem ---

% \chapter{Wstęp}
% \section{Wstęp}

Rynek walutowy stanowi jeden z najbardziej płynnych i globalnie zintegrowanych segmentów rynków finansowych.
Kurs wymiany dolara amerykańskiego względem złotego polskiego (USD/PLN) ma bezpośrednie znaczenie dla polskiej gospodarki:
wpływa na koszty importu, wartość zobowiązań zadłużonych w walucie obcej oraz poziom inflacji importowanej.
Instytucje finansowe, przedsiębiorstwa i inwestorzy indywidualni poszukują narzędzi wspierających podejmowanie decyzji na tym rynku,
jednak prognozowanie kursów walutowych pozostaje zadaniem wyjątkowo trudnym.

Szeregi czasowe notowań finansowych charakteryzują się wysokim poziomem szumu, nieliniowością zależności
oraz okresową zmianą reżimów rynkowych. Klasyczne modele ekonometryczne, oparte na założeniach stacjonarności
i liniowości, często nie wystarczają do uchwycenia złożonych wzorców zachowania cen.
W ostatnich latach coraz większą popularnością cieszą się metody uczenia głębokiego, w szczególności sieci rekurencyjne
typu LSTM (Long Short-Term Memory) oraz architektury oparte na mechanizmie uwagi (attention), takie jak Transformer.
Pozwalają one na modelowanie wielowymiarowych sekwencji czasowych i łączenie informacji z różnych horyzontów temporalnych.

\section*{Problem badawczy}

Niniejsza praca podejmuje problem prognozowania kierunku zmiany kursu USD/PLN w horyzoncie 14~dni roboczych.
Główne pytanie badawcze brzmi: \emph{czy na podstawie 30-dniowej historii cech technicznych
i makroekonomicznych można skutecznie przewidzieć, czy kurs wzrośnie, czy spadnie w ciągu kolejnych 14~dni?}

Punktem wyjścia jest skumulowany zwrot procentowy kursu zamknięcia w horyzoncie prognozy.
Ze względu na praktyczną interpretację wyników (sygnały kupna i sprzedaży) zadanie sprowadzono
do klasyfikacji binarnej: klasa~1 oznacza wzrost kursu, klasa~0 oznacza spadek.
Etykieta dla dnia $t$ definiowana jest jako:
\begin{equation}
y_t =
\begin{cases}
1 & \text{gdy } \dfrac{P_{t+14}}{P_t} - 1 > 0, \\[6pt]
0 & \text{gdy } \dfrac{P_{t+14}}{P_t} - 1 \le 0.
\end{cases}
\label{eq:target}
\end{equation}
gdzie $P_t$ oznacza kurs zamknięcia pary USD/PLN w dniu~$t$.

Skumulowany zwrot w procentach wyraża się natomiast wzorem:
\begin{equation}
    r_t = \left(\frac{P_{t+14}}{P_t} - 1\right) \cdot 100\,\%.
    \label{eq:return}
\end{equation}

\section*{Cel pracy}

Celem pracy jest zaprojektowanie, implementacja i ocena hybrydowego modelu głębokiego łączącego warstwę LSTM
z blokiem Transformer w zadaniu klasyfikacji kierunku kursu USD/PLN.
W ramach realizacji celu:
\begin{itemize}
    \item przygotowano wielowymiarowy zestaw cech obejmujący wskaźniki analizy technicznej
          oraz stopy procentowe banków centralnych (Fed, EBC, NBP),
    \item zbudowano i wytrenowano model na historycznych danych dziennych,
    \item oceniono skuteczność predykcji za pomocą metryk klasyfikacji
          (accuracy, precision, recall, F1) oraz porównano wyniki z prostymi benchmarkami
          (strategia ,,zawsze wzrost'' i ,,zawsze spadek'').
\end{itemize}

\section*{Opis metody}

\subsection*{Dane wejściowe}

Analizie poddano dzienne notowania z pliku \texttt{dane\_ekonomiczne.csv}, którego źródłem jest Yahoo Finance oraz World Bank Data.
Dane obejmują m.in.\ ceny OHLC pary USD/PLN, skorelowane pary walutowe (EUR/PLN, EUR/USD)
oraz stopy procentowe banków centralnych Stanów Zjednoczonych, strefy euro i Polski.
Z zestawu usunięto obserwacje przypadające na weekendy, aby zachować ciągłość danych sesyjnych.
Przed przekazaniem do modelu wszystkie cechy numeryczne standaryzowano za pomocą \texttt{StandardScaler},
dopasowanego wyłącznie na zbiorze treningowym.

\subsection*{Inżynieria cech}

Dla każdego dnia obliczono łącznie 24~zmienne wejściowe, w tym:
\begin{itemize}
    \item dzienne stopy zwrotu kursów zamknięcia (USD/PLN, EUR/PLN, EUR/USD),
    \item wskaźniki momentum za okresy 3, 5 i 10~dni,
    \item RSI (14~okresów), ATR (14~okresów), pozycję i szerokość pasm Bollingera (okno 20~dni),
    \item wskaźniki MACD wraz z linią sygnałową i histogramem,
    \item zmienność kroczącą, $z$-score (okno 20~dni), dzienny zakres High/Low oraz stosunek High/Low,
    \item stopy procentowe: US\_FED, EA\_ECB, PL\_NBP,
    \item opóźnione kierunki zmian ceny (lagi 1, 2, 3 i 5~dni).
\end{itemize}

\subsection*{Architektura modelu}

Model przyjmuje na wejściu sekwencję $X_t \in \mathbb{R}^{30 \times 24}$,
tj.\ 30~kolejnych dni obserwacji po 24~cechy (okno przesuwne).
Architektura hybrydowa składa się z następujących etapów:
\begin{enumerate}
    \item warstwa LSTM z 128~jednostkami ukrytymi (\texttt{return\_sequences=True}),
          wychwytująca lokalne zależności czasowe w sekwencji,
    \item warstwa \texttt{Dropout} ($p = 0{,}3$) redukująca przeuczenie,
    \item kodowanie pozycyjne (\emph{Positional Encoding}) nadające informację o położeniu obserwacji w oknie,
    \item blok Transformer z 8~głowicami uwagi (\emph{Multi-Head Attention}),
          $d_{\text{model}} = 128$ i warstwą sprzężeń zwrotnych o wymiarze 256,
          modelujący długoterminowe zależności w obrębie 30-dniowego okna,
    \item agregacja \texttt{GlobalAveragePooling1D}, po której następują warstwy gęste
          (64 i 16~neuronów, aktywacja ReLU, regularyzacja L2) oraz warstwa wyjściowa
          z aktywacją sigmoid (jedna neurona, klasyfikacja binarna).
\end{enumerate}

Połączenie LSTM i Transformera uzasadnia fakt, że sieć rekurencyjna efektywnie koduje krótkoterminowe wzorce cenowe,
natomiast mechanizm self-attention pozwala ważyć znaczenie poszczególnych dni w historycznym oknie,
co jest istotne przy analizie wielotygodniowych trendów.

\subsection*{Trenowanie i ewaluacja}

Dane podzielono chronologicznie: 80\,\% stanowi zbiór treningowy, 20\,\%~--- testowy.
Z części treningowej wydzielono 10\,\% obserwacji na walidację.
Ze względu na naturę szeregu czasowego podczas uczenia nie stosowano losowego mieszania próbek (\texttt{shuffle=False}).

Model trenowano algorytmem Adam z początkową szybkością uczenia $5 \times 10^{-5}$,
funkcją straty \texttt{binary\_crossentropy} oraz wagami klas (\texttt{compute\_class\_weight})
kompensującymi niezbalansowanie etykiet.
Zastosowano wczesne zatrzymanie (\emph{early stopping}) monitorujące dokładność walidacyjną
oraz redukcję szybkości uczenia przy stagnacji straty.
Predykcje ciągłe z warstwy sigmoid konwertowano na klasy za pomocą progu $0{,}5$.

Jakość modelu oceniano metrykami: accuracy, precision, recall, F1 oraz macierzą pomyłek.
Wyniki porównywano z naiwnymi strategiami bazowymi przewidującymi zawsze wzrost lub zawsze spadek kursu.

\section*{Zakres i ograniczenia}

Praca koncentruje się wyłącznie na prognozie \emph{kierunku} zmiany ceny, a nie na estymacji jej poziomu
ani wielkości zwrotu. Nie uwzględniono kosztów transakcyjnych, spreadu bid-ask ani dźwigni finansowej,
co ogranicza bezpośrednią aplikowalność wyników w realnym handlu algorytmicznym.

Model wytrenowano na danych historycznych z okresu od 2020~roku; jego skuteczność w przyszłości
nie jest gwarantowana ze względu na ryzyko przeuczenia oraz możliwość zmiany reżimu rynkowego
(np.\ interwencje banków centralnych, kryzysy geopolityczne). Wyniki należy traktować jako eksperyment badawczy,
a nie rekomendację inwestycyjną.

\section*{Struktura pracy}

% --- Dostosuj nazwy rozdziałów do własnego spisu treści ---

Praca składa się z następujących rozdziałów:
\begin{enumerate}
    \item Wstęp --- motywacja, problem badawczy, cel i zakres pracy.
    \item Przegląd literatury i podstawy teoretyczne --- szeregi czasowe finansowe,
          sieci LSTM, architektura Transformer oraz wybrane wskaźniki analizy technicznej.
    \item Przygotowanie danych i inżynieria cech --- opis źródeł danych, preprocessingu
          i konstrukcji 24~zmiennych wejściowych.
    \item Architektura i trenowanie modelu --- szczegółowy opis hybrydy LSTM\,+\,Transformer,
          hiperparametrów i procedury walidacji.
    \item Wyniki i analiza --- metryki klasyfikacji, macierz pomyłek, porównanie z benchmarkami
          oraz wizualizacja sygnałów kupna i sprzedaży.
    \item Podsumowanie i wnioski --- ocena osiągniętych rezultatów i kierunki dalszych badań.
\end{enumerate}

% =============================================================================

% --- Odkomentuj odpowiednią linię lub umieść treść poniżej pod własnym nagłówkiem ---

% \chapter{Wstęp}
% \section{Wstęp}

Rynek walutowy stanowi jeden z najbardziej płynnych i globalnie zintegrowanych segmentów rynków finansowych.
Kurs wymiany dolara amerykańskiego względem złotego polskiego (USD/PLN) ma bezpośrednie znaczenie dla polskiej gospodarki:
wpływa na koszty importu, wartość zobowiązań zadłużonych w walucie obcej oraz poziom inflacji importowanej.
Instytucje finansowe, przedsiębiorstwa i inwestorzy indywidualni poszukują narzędzi wspierających podejmowanie decyzji na tym rynku,
jednak prognozowanie kursów walutowych pozostaje zadaniem wyjątkowo trudnym.

Szeregi czasowe notowań finansowych charakteryzują się wysokim poziomem szumu, nieliniowością zależności
oraz okresową zmianą reżimów rynkowych. Klasyczne modele ekonometryczne, oparte na założeniach stacjonarności
i liniowości, często nie wystarczają do uchwycenia złożonych wzorców zachowania cen.
W ostatnich latach coraz większą popularnością cieszą się metody uczenia głębokiego, w szczególności sieci rekurencyjne
typu LSTM (Long Short-Term Memory) oraz architektury oparte na mechanizmie uwagi (attention), takie jak Transformer.
Pozwalają one na modelowanie wielowymiarowych sekwencji czasowych i łączenie informacji z różnych horyzontów temporalnych.

\section*{Problem badawczy}

Niniejsza praca podejmuje problem prognozowania kierunku zmiany kursu USD/PLN w horyzoncie 14~dni roboczych.
Główne pytanie badawcze brzmi: \emph{czy na podstawie 30-dniowej historii cech technicznych
i makroekonomicznych można skutecznie przewidzieć, czy kurs wzrośnie, czy spadnie w ciągu kolejnych 14~dni?}

Punktem wyjścia jest skumulowany zwrot procentowy kursu zamknięcia w horyzoncie prognozy.
Ze względu na praktyczną interpretację wyników (sygnały kupna i sprzedaży) zadanie sprowadzono
do klasyfikacji binarnej: klasa~1 oznacza wzrost kursu, klasa~0 oznacza spadek.
Etykieta dla dnia $t$ definiowana jest jako:
\begin{equation}
y_t =
\begin{cases}
1 & \text{gdy } \dfrac{P_{t+14}}{P_t} - 1 > 0, \\[6pt]
0 & \text{gdy } \dfrac{P_{t+14}}{P_t} - 1 \le 0.
\end{cases}
\label{eq:target}
\end{equation}
gdzie $P_t$ oznacza kurs zamknięcia pary USD/PLN w dniu~$t$.

Skumulowany zwrot w procentach wyraża się natomiast wzorem:
\begin{equation}
    r_t = \left(\frac{P_{t+14}}{P_t} - 1\right) \cdot 100\,\%.
    \label{eq:return}
\end{equation}

\section*{Cel pracy}

Celem pracy jest zaprojektowanie, implementacja i ocena hybrydowego modelu głębokiego łączącego warstwę LSTM
z blokiem Transformer w zadaniu klasyfikacji kierunku kursu USD/PLN.
W ramach realizacji celu:
\begin{itemize}
    \item przygotowano wielowymiarowy zestaw cech obejmujący wskaźniki analizy technicznej
          oraz stopy procentowe banków centralnych (Fed, EBC, NBP),
    \item zbudowano i wytrenowano model na historycznych danych dziennych,
    \item oceniono skuteczność predykcji za pomocą metryk klasyfikacji
          (accuracy, precision, recall, F1) oraz porównano wyniki z prostymi benchmarkami
          (strategia ,,zawsze wzrost'' i ,,zawsze spadek'').
\end{itemize}

\section*{Opis metody}

\subsection*{Dane wejściowe}

Analizie poddano dzienne notowania z pliku \texttt{dane\_ekonomiczne.csv}, którego źródłem jest Yahoo Finance oraz World Bank Data.
Dane obejmują m.in.\ ceny OHLC pary USD/PLN, skorelowane pary walutowe (EUR/PLN, EUR/USD)
oraz stopy procentowe banków centralnych Stanów Zjednoczonych, strefy euro i Polski.
Z zestawu usunięto obserwacje przypadające na weekendy, aby zachować ciągłość danych sesyjnych.
Przed przekazaniem do modelu wszystkie cechy numeryczne standaryzowano za pomocą \texttt{StandardScaler},
dopasowanego wyłącznie na zbiorze treningowym.

\subsection*{Inżynieria cech}

Dla każdego dnia obliczono łącznie 24~zmienne wejściowe, w tym:
\begin{itemize}
    \item dzienne stopy zwrotu kursów zamknięcia (USD/PLN, EUR/PLN, EUR/USD),
    \item wskaźniki momentum za okresy 3, 5 i 10~dni,
    \item RSI (14~okresów), ATR (14~okresów), pozycję i szerokość pasm Bollingera (okno 20~dni),
    \item wskaźniki MACD wraz z linią sygnałową i histogramem,
    \item zmienność kroczącą, $z$-score (okno 20~dni), dzienny zakres High/Low oraz stosunek High/Low,
    \item stopy procentowe: US\_FED, EA\_ECB, PL\_NBP,
    \item opóźnione kierunki zmian ceny (lagi 1, 2, 3 i 5~dni).
\end{itemize}

\subsection*{Architektura modelu}

Model przyjmuje na wejściu sekwencję $X_t \in \mathbb{R}^{30 \times 24}$,
tj.\ 30~kolejnych dni obserwacji po 24~cechy (okno przesuwne).
Architektura hybrydowa składa się z następujących etapów:
\begin{enumerate}
    \item warstwa LSTM z 128~jednostkami ukrytymi (\texttt{return\_sequences=True}),
          wychwytująca lokalne zależności czasowe w sekwencji,
    \item warstwa \texttt{Dropout} ($p = 0{,}3$) redukująca przeuczenie,
    \item kodowanie pozycyjne (\emph{Positional Encoding}) nadające informację o położeniu obserwacji w oknie,
    \item blok Transformer z 8~głowicami uwagi (\emph{Multi-Head Attention}),
          $d_{\text{model}} = 128$ i warstwą sprzężeń zwrotnych o wymiarze 256,
          modelujący długoterminowe zależności w obrębie 30-dniowego okna,
    \item agregacja \texttt{GlobalAveragePooling1D}, po której następują warstwy gęste
          (64 i 16~neuronów, aktywacja ReLU, regularyzacja L2) oraz warstwa wyjściowa
          z aktywacją sigmoid (jedna neurona, klasyfikacja binarna).
\end{enumerate}

Połączenie LSTM i Transformera uzasadnia fakt, że sieć rekurencyjna efektywnie koduje krótkoterminowe wzorce cenowe,
natomiast mechanizm self-attention pozwala ważyć znaczenie poszczególnych dni w historycznym oknie,
co jest istotne przy analizie wielotygodniowych trendów.

\subsection*{Trenowanie i ewaluacja}

Dane podzielono chronologicznie: 80\,\% stanowi zbiór treningowy, 20\,\%~--- testowy.
Z części treningowej wydzielono 10\,\% obserwacji na walidację.
Ze względu na naturę szeregu czasowego podczas uczenia nie stosowano losowego mieszania próbek (\texttt{shuffle=False}).

Model trenowano algorytmem Adam z początkową szybkością uczenia $5 \times 10^{-5}$,
funkcją straty \texttt{binary\_crossentropy} oraz wagami klas (\texttt{compute\_class\_weight})
kompensującymi niezbalansowanie etykiet.
Zastosowano wczesne zatrzymanie (\emph{early stopping}) monitorujące dokładność walidacyjną
oraz redukcję szybkości uczenia przy stagnacji straty.
Predykcje ciągłe z warstwy sigmoid konwertowano na klasy za pomocą progu $0{,}5$.

Jakość modelu oceniano metrykami: accuracy, precision, recall, F1 oraz macierzą pomyłek.
Wyniki porównywano z naiwnymi strategiami bazowymi przewidującymi zawsze wzrost lub zawsze spadek kursu.

\section*{Zakres i ograniczenia}

Praca koncentruje się wyłącznie na prognozie \emph{kierunku} zmiany ceny, a nie na estymacji jej poziomu
ani wielkości zwrotu. Nie uwzględniono kosztów transakcyjnych, spreadu bid-ask ani dźwigni finansowej,
co ogranicza bezpośrednią aplikowalność wyników w realnym handlu algorytmicznym.

Model wytrenowano na danych historycznych z okresu od 2020~roku; jego skuteczność w przyszłości
nie jest gwarantowana ze względu na ryzyko przeuczenia oraz możliwość zmiany reżimu rynkowego
(np.\ interwencje banków centralnych, kryzysy geopolityczne). Wyniki należy traktować jako eksperyment badawczy,
a nie rekomendację inwestycyjną.

\section*{Struktura pracy}

% --- Dostosuj nazwy rozdziałów do własnego spisu treści ---

Praca składa się z następujących rozdziałów:
\begin{enumerate}
    \item Wstęp --- motywacja, problem badawczy, cel i zakres pracy.
    \item Przegląd literatury i podstawy teoretyczne --- szeregi czasowe finansowe,
          sieci LSTM, architektura Transformer oraz wybrane wskaźniki analizy technicznej.
    \item Przygotowanie danych i inżynieria cech --- opis źródeł danych, preprocessingu
          i konstrukcji 24~zmiennych wejściowych.
    \item Architektura i trenowanie modelu --- szczegółowy opis hybrydy LSTM\,+\,Transformer,
          hiperparametrów i procedury walidacji.
    \item Wyniki i analiza --- metryki klasyfikacji, macierz pomyłek, porównanie z benchmarkami
          oraz wizualizacja sygnałów kupna i sprzedaży.
    \item Podsumowanie i wnioski --- ocena osiągniętych rezultatów i kierunki dalszych badań.
\end{enumerate}

% =============================================================================

% --- Odkomentuj odpowiednią linię lub umieść treść poniżej pod własnym nagłówkiem ---

% \chapter{Wstęp}
% \section{Wstęp}

Rynek walutowy stanowi jeden z najbardziej płynnych i globalnie zintegrowanych segmentów rynków finansowych.
Kurs wymiany dolara amerykańskiego względem złotego polskiego (USD/PLN) ma bezpośrednie znaczenie dla polskiej gospodarki:
wpływa na koszty importu, wartość zobowiązań zadłużonych w walucie obcej oraz poziom inflacji importowanej.
Instytucje finansowe, przedsiębiorstwa i inwestorzy indywidualni poszukują narzędzi wspierających podejmowanie decyzji na tym rynku,
jednak prognozowanie kursów walutowych pozostaje zadaniem wyjątkowo trudnym.

Szeregi czasowe notowań finansowych charakteryzują się wysokim poziomem szumu, nieliniowością zależności
oraz okresową zmianą reżimów rynkowych. Klasyczne modele ekonometryczne, oparte na założeniach stacjonarności
i liniowości, często nie wystarczają do uchwycenia złożonych wzorców zachowania cen.
W ostatnich latach coraz większą popularnością cieszą się metody uczenia głębokiego, w szczególności sieci rekurencyjne
typu LSTM (Long Short-Term Memory) oraz architektury oparte na mechanizmie uwagi (attention), takie jak Transformer.
Pozwalają one na modelowanie wielowymiarowych sekwencji czasowych i łączenie informacji z różnych horyzontów temporalnych.

\section*{Problem badawczy}

Niniejsza praca podejmuje problem prognozowania kierunku zmiany kursu USD/PLN w horyzoncie 14~dni roboczych.
Główne pytanie badawcze brzmi: \emph{czy na podstawie 30-dniowej historii cech technicznych
i makroekonomicznych można skutecznie przewidzieć, czy kurs wzrośnie, czy spadnie w ciągu kolejnych 14~dni?}

Punktem wyjścia jest skumulowany zwrot procentowy kursu zamknięcia w horyzoncie prognozy.
Ze względu na praktyczną interpretację wyników (sygnały kupna i sprzedaży) zadanie sprowadzono
do klasyfikacji binarnej: klasa~1 oznacza wzrost kursu, klasa~0 oznacza spadek.
Etykieta dla dnia $t$ definiowana jest jako:
\begin{equation}
y_t =
\begin{cases}
1 & \text{gdy } \dfrac{P_{t+14}}{P_t} - 1 > 0, \\[6pt]
0 & \text{gdy } \dfrac{P_{t+14}}{P_t} - 1 \le 0.
\end{cases}
\label{eq:target}
\end{equation}
gdzie $P_t$ oznacza kurs zamknięcia pary USD/PLN w dniu~$t$.

Skumulowany zwrot w procentach wyraża się natomiast wzorem:
\begin{equation}
    r_t = \left(\frac{P_{t+14}}{P_t} - 1\right) \cdot 100\,\%.
    \label{eq:return}
\end{equation}

\section*{Cel pracy}

Celem pracy jest zaprojektowanie, implementacja i ocena hybrydowego modelu głębokiego łączącego warstwę LSTM
z blokiem Transformer w zadaniu klasyfikacji kierunku kursu USD/PLN.
W ramach realizacji celu:
\begin{itemize}
    \item przygotowano wielowymiarowy zestaw cech obejmujący wskaźniki analizy technicznej
          oraz stopy procentowe banków centralnych (Fed, EBC, NBP),
    \item zbudowano i wytrenowano model na historycznych danych dziennych,
    \item oceniono skuteczność predykcji za pomocą metryk klasyfikacji
          (accuracy, precision, recall, F1) oraz porównano wyniki z prostymi benchmarkami
          (strategia ,,zawsze wzrost'' i ,,zawsze spadek'').
\end{itemize}

\section*{Opis metody}

\subsection*{Dane wejściowe}

Analizie poddano dzienne notowania z pliku \texttt{dane\_ekonomiczne.csv}, którego źródłem jest Yahoo Finance oraz World Bank Data.
Dane obejmują m.in.\ ceny OHLC pary USD/PLN, skorelowane pary walutowe (EUR/PLN, EUR/USD)
oraz stopy procentowe banków centralnych Stanów Zjednoczonych, strefy euro i Polski.
Z zestawu usunięto obserwacje przypadające na weekendy, aby zachować ciągłość danych sesyjnych.
Przed przekazaniem do modelu wszystkie cechy numeryczne standaryzowano za pomocą \texttt{StandardScaler},
dopasowanego wyłącznie na zbiorze treningowym.

\subsection*{Inżynieria cech}

Dla każdego dnia obliczono łącznie 24~zmienne wejściowe, w tym:
\begin{itemize}
    \item dzienne stopy zwrotu kursów zamknięcia (USD/PLN, EUR/PLN, EUR/USD),
    \item wskaźniki momentum za okresy 3, 5 i 10~dni,
    \item RSI (14~okresów), ATR (14~okresów), pozycję i szerokość pasm Bollingera (okno 20~dni),
    \item wskaźniki MACD wraz z linią sygnałową i histogramem,
    \item zmienność kroczącą, $z$-score (okno 20~dni), dzienny zakres High/Low oraz stosunek High/Low,
    \item stopy procentowe: US\_FED, EA\_ECB, PL\_NBP,
    \item opóźnione kierunki zmian ceny (lagi 1, 2, 3 i 5~dni).
\end{itemize}

\subsection*{Architektura modelu}

Model przyjmuje na wejściu sekwencję $X_t \in \mathbb{R}^{30 \times 24}$,
tj.\ 30~kolejnych dni obserwacji po 24~cechy (okno przesuwne).
Architektura hybrydowa składa się z następujących etapów:
\begin{enumerate}
    \item warstwa LSTM z 128~jednostkami ukrytymi (\texttt{return\_sequences=True}),
          wychwytująca lokalne zależności czasowe w sekwencji,
    \item warstwa \texttt{Dropout} ($p = 0{,}3$) redukująca przeuczenie,
    \item kodowanie pozycyjne (\emph{Positional Encoding}) nadające informację o położeniu obserwacji w oknie,
    \item blok Transformer z 8~głowicami uwagi (\emph{Multi-Head Attention}),
          $d_{\text{model}} = 128$ i warstwą sprzężeń zwrotnych o wymiarze 256,
          modelujący długoterminowe zależności w obrębie 30-dniowego okna,
    \item agregacja \texttt{GlobalAveragePooling1D}, po której następują warstwy gęste
          (64 i 16~neuronów, aktywacja ReLU, regularyzacja L2) oraz warstwa wyjściowa
          z aktywacją sigmoid (jedna neurona, klasyfikacja binarna).
\end{enumerate}

Połączenie LSTM i Transformera uzasadnia fakt, że sieć rekurencyjna efektywnie koduje krótkoterminowe wzorce cenowe,
natomiast mechanizm self-attention pozwala ważyć znaczenie poszczególnych dni w historycznym oknie,
co jest istotne przy analizie wielotygodniowych trendów.

\subsection*{Trenowanie i ewaluacja}

Dane podzielono chronologicznie: 80\,\% stanowi zbiór treningowy, 20\,\%~--- testowy.
Z części treningowej wydzielono 10\,\% obserwacji na walidację.
Ze względu na naturę szeregu czasowego podczas uczenia nie stosowano losowego mieszania próbek (\texttt{shuffle=False}).

Model trenowano algorytmem Adam z początkową szybkością uczenia $5 \times 10^{-5}$,
funkcją straty \texttt{binary\_crossentropy} oraz wagami klas (\texttt{compute\_class\_weight})
kompensującymi niezbalansowanie etykiet.
Zastosowano wczesne zatrzymanie (\emph{early stopping}) monitorujące dokładność walidacyjną
oraz redukcję szybkości uczenia przy stagnacji straty.
Predykcje ciągłe z warstwy sigmoid konwertowano na klasy za pomocą progu $0{,}5$.

Jakość modelu oceniano metrykami: accuracy, precision, recall, F1 oraz macierzą pomyłek.
Wyniki porównywano z naiwnymi strategiami bazowymi przewidującymi zawsze wzrost lub zawsze spadek kursu.

\section*{Zakres i ograniczenia}

Praca koncentruje się wyłącznie na prognozie \emph{kierunku} zmiany ceny, a nie na estymacji jej poziomu
ani wielkości zwrotu. Nie uwzględniono kosztów transakcyjnych, spreadu bid-ask ani dźwigni finansowej,
co ogranicza bezpośrednią aplikowalność wyników w realnym handlu algorytmicznym.

Model wytrenowano na danych historycznych z okresu od 2020~roku; jego skuteczność w przyszłości
nie jest gwarantowana ze względu na ryzyko przeuczenia oraz możliwość zmiany reżimu rynkowego
(np.\ interwencje banków centralnych, kryzysy geopolityczne). Wyniki należy traktować jako eksperyment badawczy,
a nie rekomendację inwestycyjną.

\section*{Struktura pracy}

% --- Dostosuj nazwy rozdziałów do własnego spisu treści ---

Praca składa się z następujących rozdziałów:
\begin{enumerate}
    \item Wstęp --- motywacja, problem badawczy, cel i zakres pracy.
    \item Przegląd literatury i podstawy teoretyczne --- szeregi czasowe finansowe,
          sieci LSTM, architektura Transformer oraz wybrane wskaźniki analizy technicznej.
    \item Przygotowanie danych i inżynieria cech --- opis źródeł danych, preprocessingu
          i konstrukcji 24~zmiennych wejściowych.
    \item Architektura i trenowanie modelu --- szczegółowy opis hybrydy LSTM\,+\,Transformer,
          hiperparametrów i procedury walidacji.
    \item Wyniki i analiza --- metryki klasyfikacji, macierz pomyłek, porównanie z benchmarkami
          oraz wizualizacja sygnałów kupna i sprzedaży.
    \item Podsumowanie i wnioski --- ocena osiągniętych rezultatów i kierunki dalszych badań.
\end{enumerate}
